Source : Eurostat, comptes annuels des administrations publiques (gov\_10a\_main, gov\_10a\_taxag, gov\_10a\_exp), PIB (nama\_10\_gdp) ; Insee, dette trimestrielle de Maastricht par sous-secteur (contributions consolidées, 4e trimestre). Calculs de l'auteur (dette/dynamique.py).\par Lecture : en 2020, le ratio de dette a augmenté de 16,8 points de PIB : 7,7 de déficit primaire, 4,8 dus à la baisse du PIB nominal (l'effet de la croissance devient positif quand le PIB recule), 1,3 d'intérêts et 3,0 d'ajustements stock-flux (trésorerie constituée par l'État). Les intérêts ajoutent chaque année de 1 à 4 points ; la croissance en retire en général autant.
